\documentclass[12pt,a4paper]{article}
\usepackage[utf8]{inputenc}
\usepackage[T1]{fontenc}
\usepackage{geometry}
\usepackage{listings}
\usepackage{xcolor}
\usepackage{hyperref}
\usepackage{titlesec}

\geometry{margin=0.8in}

\definecolor{codegreen}{rgb}{0,0.6,0}
\definecolor{codegray}{rgb}{0.5,0.5,0.5}
\definecolor{codepurple}{rgb}{0.58,0,0.82}
\definecolor{backcolour}{rgb}{0.97,0.97,0.95}

\lstdefinestyle{mystyle}{
    backgroundcolor=\color{backcolour},   
    commentstyle=\color{codegreen},
    keywordstyle=\color{blue},
    numberstyle=\tiny\color{codegray},
    stringstyle=\color{codepurple},
    basicstyle=\ttfamily\small,
    breakatwhitespace=false,         
    breaklines=true,                 
    captionpos=b,                    
    keepspaces=true,                 
    numbers=left,                    
    numbersep=5pt,                  
    showspaces=false,                
    showstringspaces=false,
    showtabs=false,                  
    tabsize=2,
    frame=single
}

\lstset{style=mystyle}

\title{\textbf{Detailed Technical Report: JEE Backend Architecture}}
\author{Student Name: Mounib \\ Course: JEE L2-S2}
\date{\today}

\begin{document}

\maketitle
\tableofcontents
\newpage

\section{Project Overview}
This project is a comprehensive Student Management System built on the Jakarta EE 10 platform. The backend is designed as a \textbf{RESTful API} that provides data services to a JavaFX client. It leverages \textbf{EJBs} for business logic, \textbf{JPA} for database persistence, and \textbf{JAX-RS} for network communication.

\section{Core Configuration Files}

\subsection{persistence.xml (The Database Heart)}
Located in \texttt{src/main/resources/META-INF/}, this file configures the Java Persistence API.
\begin{itemize}
    \item \textbf{jta-data-source}: Points to \texttt{java:jboss/datasources/ExampleDS}, which is the built-in H2 database in WildFly.
    \item \textbf{drop-and-create}: Instructs Hibernate to reset the database structure on every deploy.
    \item \textbf{sql-load-script-source}: Automatically runs \texttt{import.sql} to seed the database with initial data.
\end{itemize}

\subsection{HelloApplication.java (The API Entry)}
\begin{lstlisting}[language=Java]
@ApplicationPath("/api")
public class HelloApplication extends Application { }
\end{lstlisting}
\textbf{Role}: This class activates JAX-RS. The \texttt{@ApplicationPath("/api")} annotation defines the base URL for all your REST endpoints.

\section{The Entity Layer (Data Models)}

\subsection{User.java \& Student.java}
These are \textbf{POJOs} (Plain Old Java Objects) annotated with \texttt{@Entity}.
\begin{itemize}
    \item \texttt{@Id} \& \texttt{@GeneratedValue}: Automatically manages the primary key (ID).
    \item \texttt{@Table(name="...")}: Maps the Java class to a specific SQL table name.
    \item \texttt{Serializable}: Allows the object to be converted into a stream of bytes for network transmission or persistence.
\end{itemize}

\section{The Repository Layer (Database Operations)}

\subsection{StudentRepository.java}
This is a \textbf{Stateless EJB} that handles SQL operations through the \texttt{EntityManager}.
\begin{itemize}
    \item \textbf{findAll()}: Executes a \textit{JPQL} query (\texttt{SELECT s FROM Student s}) to retrieve all records.
    \item \textbf{save(Student s)}: Uses \texttt{persist()} if it's a new student, or \texttt{merge()} if it's an update. This is the core of the **Create/Update** functionality.
    \item \textbf{delete(Long id)}: Finds the object in the database and calls \texttt{remove()}.
\end{itemize}

\section{The Resource Layer (REST API Endpoints)}

\subsection{AuthResource.java (Authentication Logic)}
\textbf{Function: login(User request)}
\begin{enumerate}
    \item Receives JSON from the JavaFX client.
    \item Calls \texttt{findByUsername()} to check if the user exists.
    \item \textbf{Security Check}: Calls \texttt{passwordService.verifyPassword()}.
    \item \textbf{Response}: Returns a JSON object containing a session token and the user's role.
\end{enumerate}

\subsection{StudentResource.java (Data Exchange)}
This class connects the network to the database logic.
\begin{itemize}
    \item \texttt{@GET}: Returns the list of students. Jakarta EE automatically converts the Java List into a \textbf{JSON Array}.
    \item \texttt{@POST}: Receives a new student as JSON, converts it to a Java Object, and saves it.
    \item \texttt{@DELETE}: Receives an ID from the URL path and removes the record.
\end{itemize}

\section{Security Implementation}

\subsection{PasswordService.java}
Uses the \texttt{Pbkdf2PasswordHash} interface from Jakarta Security.
\begin{itemize}
    \item \textbf{Hashing}: Transforms a plain password into a secure string using 2048 iterations and a 32-byte salt.
    \item \textbf{Verification}: Compares the input password against the stored hash safely.
\end{itemize}

\section{The REST-to-Client Connection}
The communication is handled via the \textbf{JSON (JavaScript Object Notation)} format.
\begin{itemize}
    \item \textbf{Serialization}: The Backend transforms Java objects into JSON text.
    \item \textbf{Deserialization}: The Frontend (JavaFX) receives this text and parses it back into Java models.
    \item \textbf{HTTP Methods}: 
    \begin{itemize}
        \item \texttt{GET} for reading data.
        \item \texttt{POST} for creating data.
        \item \texttt{DELETE} for removing data.
    \end{itemize}
\end{itemize}

\section{Conclusion}
This architecture ensures that the Backend is \textbf{Stateless} and \textbf{Independent}. This means you could theoretically replace your JavaFX frontend with a Web or Mobile app, and the Backend would continue to work perfectly without a single line of code being changed.

\end{document}
